\documentclass[12pt,a4paper]{article}

\usepackage{graphicx}
\usepackage{geometry}
\usepackage{fancyhdr}
\usepackage{xcolor}
\usepackage{array}
\usepackage{titlesec}

\geometry{margin=1in}

% Blue section headings
\titleformat{\section}
{\color{blue}\normalfont\Large\bfseries}
{\thesection}{1em}{}

\begin{document}

% Header
\begin{minipage}{0.3\textwidth}
    \includegraphics[width=\linewidth]{LOGO.PNG}
\end{minipage}
\begin{minipage}{0.8\textwidth}
    \hfill
    \begin{tabular}{rl}
        \textbf{Name:} & Pavan Kumar KS \\
        \textbf{ID:} & COMETFWC044 \\
        \textbf{Date:} & 23 Feb 2026 \\
    \end{tabular}
\end{minipage}

\vspace{1cm}

%------------------------------------------------
\section*{QUESTION}

A 4:1 Multiplexer has select lines P and Q.  
The inputs are given as:

\[
I_0 = R,\quad
I_1 = \bar{R},\quad
I_2 = \bar{R},\quad
I_3 = R
\]

Find the Boolean expression for the output function F.

%------------------------------------------------
\section*{QUESTION ANALYSIS}

For a 4:1 MUX,

\[
F = I_0\bar{P}\bar{Q} + I_1\bar{P}Q + I_2P\bar{Q} + I_3PQ
\]

Substituting the inputs:

\[
F = R\bar{P}\bar{Q} + \bar{R}\bar{P}Q + \bar{R}P\bar{Q} + RPQ
\]

Grouping terms:

\[
F = R(\bar{P}\bar{Q} + PQ) + \bar{R}(\bar{P}Q + P\bar{Q})
\]

\[
F = R(P \oplus Q)' + \bar{R}(P \oplus Q)
\]

\[
\boxed{F = P \oplus Q \oplus R}
\]

%------------------------------------------------
\section*{TRUTH TABLE}

\begin{center}
\begin{tabular}{|c|c|c|c|}
\hline
P & Q & R & F = P $\oplus$ Q $\oplus$ R \\
\hline
0 & 0 & 0 & 0 \\
0 & 0 & 1 & 1 \\
0 & 1 & 0 & 1 \\
0 & 1 & 1 & 0 \\
1 & 0 & 0 & 1 \\
1 & 0 & 1 & 0 \\
1 & 1 & 0 & 0 \\
1 & 1 & 1 & 1 \\
\hline
\end{tabular}
\end{center}

%------------------------------------------------
\section*{HARDWARE IMPLEMENTATION}

The expression is implemented using:
Arduino UNO + 7447 BCD to 7-segment driver + Common Anode 7-segment display.

%------------------------------------------------
\section*{COMPONENTS REQUIRED}

\begin{center}
\begin{tabular}{|l|c|}
\hline
Component & Quantity \\
\hline
Arduino UNO & 1 \\
7447 IC & 1 \\
Common Anode 7-Segment Display & 1 \\
220$\Omega$ Resistors & 7 \\
Breadboard & 1 \\
Jumper Wires & As required \\
\hline
\end{tabular}
\end{center}

%------------------------------------------------
\section*{PIN CONNECTIONS}

\begin{center}
\begin{tabular}{|l|l|}
\hline
Connection & Details \\
\hline
Arduino D2 & Input P \\
Arduino D3 & Input Q \\
Arduino D4 & Input R \\
Arduino D8 & 7447 Pin 7 (A) \\
Arduino D9 & 7447 Pin 1 (B) \\
Arduino D10 & 7447 Pin 2 (C) \\
Arduino D11 & 7447 Pin 6 (D) \\
7447 Pin 16 & +5V \\
7447 Pin 8 & GND \\
7447 Outputs (a-g) & To 7-segment via 220$\Omega$ resistors \\
Common Anode & +5V \\
\hline
\end{tabular}
\end{center}

%------------------------------------------------
\section*{LOGIC DESCRIPTION}

The Arduino reads inputs P, Q and R.  
It computes the XOR operation:

\[
F = P \oplus Q \oplus R
\]

The output F is sent as BCD input to the 7447 IC.  
The 7447 decodes the BCD value and drives the 7-segment display.  

If the number of HIGH inputs is odd → display shows 1.  
If the number of HIGH inputs is even → display shows 0.

%------------------------------------------------
\section*{CODE UPLOADING STEPS}

\begin{enumerate}
\item Create a blink folder and write Makefile
\item Write the code in main.c
\item Run the command ``make''. It will compile the code and create hex file
\item Copy the .hex file to ArduinoDroid folder
\item Connect the Arduino UNO to mobile with OTG cable
\item Upload the hex file using ``Upload Precompiled'' option
\item Observe the output and verify the expression
\end{enumerate}

%------------------------------------------------
\section*{EXPERIMENTAL TRUTH TABLE}

\begin{center}
\begin{tabular}{|c|c|c|c|}
\hline
P & Q & R & Observed Output (F) \\
\hline
0 & 0 & 0 & 0 \\
0 & 0 & 1 & 1 \\
0 & 1 & 0 & 1 \\
0 & 1 & 1 & 0 \\
1 & 0 & 0 & 1 \\
1 & 0 & 1 & 0 \\
1 & 1 & 0 & 0 \\
1 & 1 & 1 & 1 \\
\hline
\end{tabular}
\end{center}
%------------------------------------------------
\section*{CIRCUIT CONNECTION}

\begin{center}
\includegraphics[width=0.8\textwidth]{pic.jpg}
\end{center}

%------------------------------------------------
\section*{CONCLUSION}

The given 4:1 multiplexer configuration simplifies to the XOR expression:

\[
F = P \oplus Q \oplus R
\]

The hardware implementation using Arduino and 7447 successfully verifies the Boolean expression through the 7-segment display output.

\end{document}
